\section{Sparse quadratic hulls and graph minors}
\label{graph:section}

The dense four-variable obstruction extends to sparse models through graph
minors. The distinction between positive and negative diagonal coordinates
is essential. A nonnegative square coefficient pairs with a diagonal
epigraph, whereas a nonpositive square coefficient pairs with a diagonal
hypograph.

Let $G=(V,E,L_+,L_-)$, where $E$ is a set of unordered pairs of distinct
vertices and $L_+,L_-\subseteq V$ are disjoint. Put $P=L_+$ and $R=V\setminus P$.
The retained coordinates are $m_i$ for $i\in V$, $Y_{ij}$ for $ij\in E$,
and $Y_{ii}$ for $i\in L_+\cup L_-$. Define the compact graph moment hull
\[
 \mathcal K(G)=\conv\{(x,(x_ix_j)_{ij\in E},
                         (x_i^2)_{i\in L_+\cup L_-}):x\in[0,1]^V\}.
\]
Let $D(G)$ add nonnegative slack on positive diagonals and nonpositive
slack on negative diagonals, with zero entries on all other retained
coordinates. The signed hull is
\begin{equation}\label{graph:signed-hull}
 \mathcal H(G)=\mathcal K(G)+D(G).
\end{equation}
Equivalently, it is the convex hull of points with exact edge products,
$Y_{ii}\geq x_i^2$ for $i\in L_+$, and $Y_{ii}\leq x_i^2$ for
$i\in L_-$. It is closed because $\mathcal K(G)$ is compact and $D(G)$ is
closed. This is the graph hull studied by \citet{Khajavirad2026SparseBoxQP}.

\subsection{A \texorpdfstring{$K_4$}{K4} minor among positive vertices}

\begin{theorem}\label{graph:k4}
If the graph $G[P]$ induced by the positive-loop vertices contains a
$K_4$ minor, then $\mathcal H(G)$ has no finite semidefinite lift.
\end{theorem}
\begin{proof}
Choose disjoint connected branch sets $C_1,\ldots,C_4\subseteq P$ witnessing
the minor. In each nonsingleton branch set choose a spanning tree, and
choose one edge $e_{ab}$ between $C_a$ and $C_b$ for every $a<b$.
The affine function
\begin{equation}\label{graph:contraction-face}
 F=\sum_{ij\text{ in the chosen trees}}
       (Y_{ii}+Y_{jj}-2Y_{ij})
\end{equation}
is nonnegative on $\mathcal H(G)$. Indeed, each summand at a generating
point is $(x_i-x_j)^2+s_i+s_j$, with $s_i,s_j\geq0$.

On the face $F=0$, every positive-weight atom is constant on each branch
set. All diagonal slack on nonsingleton branch sets vanishes, since each
vertex belongs to a tree edge. Slack at singleton branch sets may remain.
Choose a representative $r_a\in C_a$ and project this face onto
\[
 u_a=m_{r_a},\qquad Z_{aa}=Y_{r_ar_a},\qquad
 Z_{ab}=Y_{e_{ab}}\quad(a<b).
\]
The image is exactly $\mathcal Q_4+D_S$, where $D_S$ is nonnegative
diagonal slack on the singleton branch sets. For the forward inclusion,
the common branch values of each atom give a cube atom in four variables,
and only those singleton slacks survive. For the reverse inclusion, assign
the coordinate $u_a$ of any cube atom to every vertex in $C_a$, assign zero
outside the branch sets, and take exact squares and products in the source
graph. Add any desired singleton slack and take mixtures.

If $\mathcal H(G)$ had a finite lift, so would its face and this projection.
Adding the remaining nonnegative diagonal rays would then give a finite
lift of $\mathcal H_4$, contradicting Theorem~\ref{repr:threshold}.
\end{proof}

Extra edges, extra vertices, and negative-loop vertices do not affect the
proof. The branch sets must lie in $P$: a negative or absent diagonal cannot
provide the nonnegative squared-difference expression in
\eqref{graph:contraction-face}. The theorem therefore gives an obstruction
in the positive-induced graph, rather than in the full interaction graph.
It applies, for example, to every subdivision of $K_4$ with positive loops
at all subdivision vertices.

$K_4$ already prevents a finite lift, although every proper face of its
compact moment hull has one by Proposition~\ref{repr:proper-faces}.

\subsection{Minor closure for all-positive graphs}

When every vertex has a positive loop, write $\mathcal H(G)$ for this
all-positive hull.

\begin{proposition}\label{graph:minor-closure}
If $\mathcal H(G)$ has a finite semidefinite lift and $G'$ is a minor of $G$,
then $\mathcal H(G')$ has a finite semidefinite lift.
\end{proposition}
\begin{proof}
Deleting an edge or a vertex is a coordinate projection. To contract an
edge $ij$, impose the exposed-face equation
$Y_{ii}+Y_{jj}-2Y_{ij}=0$. As in
\eqref{graph:contraction-face}, each atom then satisfies $x_i=x_j$, and
the two diagonal slacks vanish. Keep one first-moment coordinate and one
diagonal coordinate for the merged vertex. For an edge from the merged
vertex to a neighbor $k$, keep either $Y_{ik}$ or $Y_{jk}$ that is present;
if both are present, they agree on this face. Keep all other coordinates.
The projection has the graph moment hull of $G/ij$, with nonnegative
diagonal slack at every vertex except the merged vertex. Add the diagonal
ray at that vertex to obtain exactly $\mathcal H(G/ij)$.
All three operations preserve finite lifts.
\end{proof}

The diagonal-ray addition is necessary. The contraction face forces the
merged diagonal to be an exact second moment; merely projecting that face
does not recover its full positive diagonal epigraph.

\subsection{Finite lifts for small positive components}

At the other end, components with at most three positive vertices admit
finite lifts. This combines binary rounding outside $P$, the established
three-variable moment lift, and consistent binary marginal tables.
We give the formulation because it makes clear which graph width controls
its size.

For each connected component $C$ of $G[P]$, let
$B_C=N(C)\cap R$. Form a graph $T$ on $R$ by keeping its original edges
and making every $B_C$ a clique. We call $T$ the \emph{torso} after the
positive components are removed.

\begin{theorem}\label{graph:small-components}
If every connected component of $G[P]$ has at most three vertices, then
$\mathcal H(G)$ has a finite semidefinite lift whose matrix blocks have
order at most four.

More precisely, suppose a tree decomposition of $T$ with bags
$\mathcal B$ and width $\tau$ is supplied. The construction uses
$\sum_{B\in\mathcal B}2^{|B|}$ nonnegative binary table entries and
\begin{equation}\label{graph:block-count}
 \sum_C |C|!\,2^{|B_C|}
\end{equation}
positive semidefinite blocks of order $|C|+1$, each constrained to be
entrywise nonnegative. Apart from the original retained coordinates, the
number of variables and affine constraints is
\[
 O\left(|\mathcal B|2^{\tau+1}
          +\sum_C |C|!\,2^{|B_C|}\right).
\]
This count does not count every nonzero coefficient in the displayed
linear maps.
\end{theorem}
\begin{proof}
First make the coordinates in $R$ binary. At an original generating point,
independently round $x_j$ to a Bernoulli variable with mean $x_j$ for every
$j\in R$, leaving coordinates in $P$ fixed. The first moments and every
edge product remain unchanged. A negative-loop coordinate can be restored
by downward slack, because the rounded square has mean $x_j\geq x_j^2$.
Positive diagonal coordinates are unchanged. Thus the original generator
lies in the convex hull of generators with binary $R$. Conversely, those
binary generators are allowed original generators. It is enough to
describe their mixtures.

For each bag $B$ of the decomposition introduce a nonnegative table
$p_B:\{0,1\}^B\to\R_+$ with total mass one. Require adjacent bags to
agree on the marginal table of their intersection. These tables admit a
joint binary law on $R$. To construct it, root the decomposition tree and
attach each child table by its conditional law given the separator state.
Separator states of zero mass never occur and may be assigned arbitrary
conditional laws. The running-intersection property ensures that the
joint law has every prescribed bag marginal. This construction is finite.
If $R$ is empty, use a single empty bag whose sole table entry is $1$;
the empty boundary assignment then has probability $1$.

Each $B_C$ is a clique in $T$, so it is contained in a bag. For every
$b\in\{0,1\}^{B_C}$ introduce the homogenized simplex lift of the component
cube $[0,1]^C$ with total mass equal to the marginal probability of $b$.
Writing $k=|C|$ and $V_\pi$ for an order-simplex vertex matrix, its variables
satisfy
\begin{align*}
 &W_{C,b,\pi}\succeq0,\qquad W_{C,b,\pi}\geq0,\\
 &\sum_\pi\mathbf1^\top W_{C,b,\pi}\mathbf1
      =\Pr(X_{B_C}=b),\\
 &m_{C,b}=\sum_\pi V_\pi W_{C,b,\pi}\mathbf1,
 \qquad Z_{C,b}=\sum_\pi V_\pi W_{C,b,\pi}V_\pi^\top.
\end{align*}
Because $k+1\leq4$, $\CP_{k+1}=\DNN_{k+1}$ proves exactness just as in
\eqref{repr:small-lift}. Zero mass forces every entry of every $W$ to be
zero, so no conditional moments can remain at a zero-probability state.

Sum $m_{C,b}$ and $Z_{C,b}$ over $b$ for component first and second moments.
For an edge $ij$ with $i\in C$ and $j\in B_C$, use
\[
 Y_{ij}=\sum_b b_j(m_{C,b})_i.
\]
Moments on $R$ come from the bag tables. At a negative-loop vertex $j$,
set $Y_{jj}=m_j-t_j$ with $t_j\geq0$; at a positive-loop vertex add
nonnegative diagonal slack to the component second moment.

Every original generator and every mixture gives feasible variables.
Conversely, begin with the joint binary law on $R$ just constructed.
Given the full binary assignment, sample each continuous component from
the law specified by its boundary assignment $b$, independently across
components. The simplex decomposition supplies these component laws.
There are no edges between distinct positive components. Hence this joint
law realizes every retained first moment, edge product, and unslacked
diagonal. Adding the stated signed slacks proves reverse inclusion.

The table count is immediate. Each pair $(C,b)$ has $k!$ simplex blocks,
each of bounded order. Separator equalities use at most
$O(|\mathcal B|2^{\tau+1})$ entries, and attachment equalities use one
equation per boundary state. This proves the stated count.
\end{proof}

For finite representability alone one may take a single bag containing
all of $R$; the formulation can then be exponentially large. A small
formulation follows when the supplied torso decomposition and the sum
in \eqref{graph:block-count} give a small total size. Original graph
treewidth alone does not control this count: completing component
boundaries to cliques can increase torso width. The theorem also concerns
quadratic graph moments; it does not describe a retained product of three
continuous coordinates.

Theorem~\ref{graph:small-components} extends the two-positive-vertex
component construction of \citet{Khajavirad2026SparseBoxQP} by using the known
three-variable simplex lift. Its sufficiency and
Theorem~\ref{graph:k4} leave an explicit gap: components with at least four
vertices and no $K_4$ minor are not classified here.

\subsection{The remaining path and star question}

Every connected graph with at least four vertices contains a four-vertex
path $P_4$ or a three-leaf star $K_{1,3}$ as a subgraph. Indeed, a spanning
tree either has a vertex of degree at least three or is a path. Together
with Proposition~\ref{graph:minor-closure} and
Theorem~\ref{graph:small-components}, this gives a precise reduction: if
both $\mathcal H(P_4)$ and $\mathcal H(K_{1,3})$ lack finite lifts, then an
all-positive graph has a finite lift exactly when every component has at
most three vertices. We do not assume either missing result.

There is also an exact reduction of the forest question to submodular
quadratics. Let $\mathcal P_G^+$ be the cone of quadratics nonnegative on
$[0,1]^V$, with square coefficients nonnegative and cross coefficients
zero outside $E$. Let
\[
 \mathcal S_-(G)=\{q\in\mathcal P_G^+:c_{ij}\leq0\text{ for }ij\in E\}.
\]
These are the submodular quadratics in the sparse cone.

\begin{proposition}\label{graph:forest-submodular}
For a forest $G$, the hull $\mathcal H(G)$ has a finite semidefinite lift
if and only if $\mathcal S_-(G)$ is a spectrahedral shadow.
\end{proposition}
\begin{proof}
Complementing a coordinate, $x_i\mapsto1-x_i$, preserves box nonnegativity
and every square coefficient. It reverses the sign of each incident cross
coefficient. On a forest the complementations can be chosen independently
along edges: root each tree, choose the root sign, and then choose each
child sign to make its parent-edge coefficient nonpositive. Therefore
every member of $\mathcal P_G^+$ is the complementation image of a member
of $\mathcal S_-(G)$. Each such image remains in $\mathcal P_G^+$, so
\[
 \mathcal P_G^+=\sum_\sigma \sigma(\mathcal S_-(G)),
\]
where the finite sum ranges over coordinate complementations, acting
linearly and invertibly on polynomial coefficients.
If $\mathcal S_-(G)$ is a shadow, the sum is a shadow. Conversely,
$\mathcal S_-(G)$ is obtained from $\mathcal P_G^+$ by linear inequalities
on the cross coefficients and therefore is a shadow whenever
$\mathcal P_G^+$ is. Lemma~\ref{repr:normalization} and closed-cone duality
identify finite representability of $\mathcal P_G^+$ with that of
$\mathcal H(G)$, proving the assertion.
\end{proof}

In particular, a finite lift for the submodular part of the dense
four-variable positive cone would yield lifts for both $P_4$ and
$K_{1,3}$ by setting the other cross coefficients to zero. The dense
nonrepresentability theorem does not decide this sign-restricted question.
Nor does failure of a particular relaxation decide it: the issue is the
existence of an arbitrary finite extended formulation.
